\documentclass[10pt,aspectratio=169]{beamer}
\input{../../_en_preambulo_beamer}% Comandos y macros estadísticas compartidas para las presentaciones Beamer en inglés (Metropolis)

% Conjuntos numéricos básicos
\newcommand{\R}{\mathbb{R}}
\newcommand{\N}{\mathbb{N}}
\newcommand{\Q}{\mathbb{Q}}
\newcommand{\Z}{\mathbb{Z}}
\newcommand{\set}[1]{\left\{ #1 \right\}}
\newcommand{\abs}[1]{\left|#1\right|}
\newcommand{\norm}[1]{\left\| #1 \right\|}

% Comandos estadísticos y probabilísticos del curso
\newcommand{\Var}{\operatorname{Var}}
\newcommand{\cov}{\operatorname{Cov}}
\newcommand{\corr}{\rho}
\newcommand{\comb}[2]{\begin{pmatrix} #1 \\ #2 \end{pmatrix}}
\newcommand{\s}{\sigma}
\newcommand{\particion}{\mathfrak{P}}
\newcommand{\card}[1]{\#\left( #1 \right)}
\newcommand{\seq}[3]{\set{#1}_{#2}^{#3}}
\newcommand{\E}{\mathbb{E}}
\newcommand{\Prob}{\mathbb{P}}

% Entornos pedagógicos y cajas en inglés académico natural (soportando tanto nombres en inglés como en español)
\newenvironment{discussion}{%
  \begin{alertblock}{Discussion Question}%
}{%
  \end{alertblock}%
}
\newenvironment{discusion}{%
  \begin{alertblock}{Discussion Question}%
}{%
  \end{alertblock}%
}

\newenvironment{reflection}{%
  \begin{alertblock}{Classroom Reflection}%
}{%
  \end{alertblock}%
}
\newenvironment{reflexion}{%
  \begin{alertblock}{Classroom Reflection}%
}{%
  \end{alertblock}%
}

\newenvironment{paradox}{%
  \begin{alertblock}{Exploratory Paradox}%
}{%
  \end{alertblock}%
}
\newenvironment{paradoja}{%
  \begin{alertblock}{Exploratory Paradox}%
}{%
  \end{alertblock}%
}

\newenvironment{motivation}{%
  \begin{exampleblock}{Motivation: Data Science in Practice}%
}{%
  \end{exampleblock}%
}
\newenvironment{motivacion}{%
  \begin{exampleblock}{Motivation: Data Science in Practice}%
}{%
  \end{exampleblock}%
}

\newenvironment{quickchallenge}{%
  \begin{block}{Quick Team Challenge}%
}{%
  \end{block}%
}
\newenvironment{retoclase}{%
  \begin{block}{Quick Team Challenge}%
}{%
  \end{block}%
}

\title[Hypothesis testing guide]{Section 07.05: A Step-by-Step Hypothesis Test Guide}\subtitle{From question to report}
\author[J. Castillo Colmenares]{Juliho Castillo Colmenares}\institute{Tecnológico de Monterrey}\date{}
\begin{document}
\begin{frame}[plain]\titlepage\end{frame}
\begin{frame}{Roadmap}\small\begin{enumerate}\item State hypotheses.\item Sample and standardize.\item Compare, decide, and communicate.\end{enumerate}\end{frame}
\begin{frame}{Source and scope}\small This guide summarizes the notes as a reproducible sequence.\begin{block}{Guiding question}Which decisions must be documented?\end{block}\end{frame}
\begin{frame}{Step 1: null hypothesis}\small The null states the reference claim about the parameter and contains $A_0$.\[H_0:A=A_0.\]\end{frame}
\begin{frame}{Step 1: alternative}\small The alternative states the departure of interest: $A<A_0$, $A>A_0$, or $A\neq A_0$. Its direction sets the tails.\end{frame}
\begin{frame}{Step 2: sample}\small Take a random sample of units or occurrences and compute the relevant estimator, such as a mean $A_m$.\begin{alertblock}{Quality} Independence and the sampling plan support interpretation.\end{alertblock}\end{frame}
\begin{frame}{Step 2: statistic}\small Summarize the difference between $A_m$ and $A_0$ using a standard error. For known $\sigma$, $Z=(A_m-A_0)/(\sigma/\sqrt n)$.\end{frame}
\begin{frame}{Step 3: normal value}\small The notes introduce the standard-normal value $Z$ and its associated probability. The null distribution determines the relevant area.\end{frame}
\begin{frame}{Step 3: levels and tails}\small Set the significance level before observing data. A one-sided alternative uses one tail; a two-sided alternative splits $\alpha$.\end{frame}
\begin{frame}{Step 4: comparison}\small Compare the probability associated with the statistic to $\alpha$.\[p\leq\alpha\Rightarrow\text{reject};\qquad p>\alpha\Rightarrow\text{do not reject}.\]\end{frame}
\begin{frame}{Step 5: decision}\small The decision is conditional on the model and selected level. Do not reject does not mean definitive acceptance of the null.\end{frame}
\begin{frame}{Step 6: communication}\small Report the hypotheses, statistic, level, $p$-value, and a conclusion in the context of the measured variable.\end{frame}
\begin{frame}{Application: restaurant}\small A restaurant claims a mean delivery time of $20$ minutes with standard deviation $3$ minutes. The claim is evaluated with a sample.\end{frame}
\begin{frame}{Application: setup}\small\[H_0:\mu=20.\]The alternative expresses the concern: $\mu>20$, $\mu<20$, or $\mu\neq20$.\end{frame}
\begin{frame}{Application: standardization}\small With sample mean $\bar x$ and size $n$,\[Z=\frac{\bar x-20}{3/\sqrt n}.\]The sign gives the direction of evidence.\end{frame}
\begin{frame}{Application: tail area}\small Obtain the probability associated with $Z$ from the standard normal and convert it to a $p$-value for the chosen alternative.\end{frame}
\begin{frame}{Application: decision}\small Compare the $p$-value with $\alpha$. State whether evidence is sufficient to question the twenty-minute claim.\end{frame}
\begin{frame}{Common errors}\small\begin{itemize}\item Swap $H_0$ and $H_a$.\item Choose the tail after computing.\item Confuse non-rejection with acceptance.\item Omit units and context.\end{itemize}\end{frame}
\begin{frame}{Checklist}\small\begin{enumerate}\item Parameter identified.\item Complete hypotheses.\item Assumptions checked.\item $\alpha$ set.\item Statistic and $p$-value reported.\end{enumerate}\end{frame}
\begin{frame}{Reporting template}\small\begin{block}{Wording} ``For $H_0:\cdots$ versus $H_a:\cdots$, we obtained $T=\cdots$ and $p=\cdots$. At $\alpha=\cdots$, we [reject/do not reject] $H_0$.''\end{block}\end{frame}
\begin{frame}{Synthesis}\small The sequence turns an applied question into an auditable, communicable statistical decision.\end{frame}
\begin{frame}{Chapter connection}\small The steps recur in $t$ tests, proportion, variance, homogeneity, and multiple-proportion tests.\end{frame}
\end{document}
